\section{Implementation and Validation Details}

\subsection{Repository Engineering Architecture}

The system is organized so that each layer can be tested independently and integrated through explicit contracts:
\begin{itemize}
  \item Core blockchain primitives expose deterministic block and transaction semantics.
  \item RPC interfaces use schema-driven request/response models, reducing ambiguity at API boundaries.
  \item AI modules are separated into encoder, policy, value, and trainer code paths to isolate optimization concerns.
  \item Quantum modules (QNN and QTransformer) are provided as pluggable components with benchmarked runtime footprints.
  \item Experiment scripts produce machine-readable outputs that are post-processed into paper tables and plots.
\end{itemize}

This decomposition follows reproducible systems guidance where experiment claims are linked to executable artifacts rather than manually curated summaries \citep{wilson2017good,stodden2016enhancing}.

\subsection{Validation Gates}

The project gate stack includes:
\begin{enumerate}
  \item static lint checks,
  \item static type checks,
  \item unit and integration tests,
  \item API contract tests,
  \item runtime health checks,
  \item manuscript artifact regeneration.
\end{enumerate}

The intent is to prevent a common anti-pattern in research software: result-generation scripts diverging from serving/runtime paths. In QAI-Chain, manuscript figures and tables are generated from tracked JSON artifacts, which are themselves emitted by executable protocol scripts.

\subsection{Latency and System Cost Profile}
\label{sec:latency-profile}

Using the latency profile shown in Figure~\ref{fig:core-latency}, QTransformer latency dominates end-to-end inference cost, with QNN as the second-largest contributor and the classical metrics encoder effectively negligible at this scale. This motivates two engineering strategies:
\begin{itemize}
  \item selective invocation of quantum modules only when predictive uncertainty exceeds threshold,
  \item asynchronous or batched execution for quantum pathways in throughput-critical deployments.
\end{itemize}

\subsection{Statistical Reporting Choices}

We report effect sizes and permutation tests to avoid over-reliance on point estimates. Non-parametric tests are particularly useful at small sample sizes and under non-Gaussian reward distributions. We include bootstrap confidence intervals to characterize uncertainty around means and to help distinguish practical from purely numerical differences \citep{efron1994introduction,good2005permutation}.

\subsection{Failure Modes Observed During Development}

A practical contribution of this work is the debugging record converted into engineering controls. We encountered and addressed:
\begin{itemize}
  \item tensor-to-NumPy conversion failures with gradient-bearing tensors,
  \item qubit-feature dimensionality mismatches in angle embedding pathways,
  \item import path drift between utility and logger modules,
  \item schema and payload mismatches in endpoint handlers.
\end{itemize}

The resulting fixes are reflected in tests and health checks, reducing recurrence probability and improving confidence in subsequent experiments.
